\subsection{向量丛截面的喷射}

12.6.1. 设 E 是以 X 为底空间的类 \(C^{r}\) 向量丛，且 \(\pi\) 是其投影。令 \(c_{0} = (\mathrm{U}, \varphi, \mathrm{F}_{0})\) 为 E 的向量图表，令 \(c_{1} = (\mathrm{U}, \psi, \mathrm{F}_{1})\) 为流形 X 的图表，且定义域同为 U。这些图表定义了一个双射 \(\theta\)，从 \(\mathrm{P}^{k}(\pi)|\mathrm{U}\) 到 \(\mathrm{J}^{k}(\psi(\mathrm{U}), \mathrm{F}_{0}) = \psi(\mathrm{U}) \times \mathrm{F}_{0} \times \mathrm{Q}^{k}(\mathrm{F}_{1}, \mathrm{F}_{0})\)（参见 12.2.2 和 12.3.1），由此得到一个向量图表

\[
d = (\mathrm{U}, \theta, \mathrm{G}), \quad \text{其中}\quad \mathrm{G} = \mathrm{F}_{0}\times\mathrm{Q}^{k}(\mathrm{F}_{1},\mathrm{F}_{0}) = \prod_{m=0}^{k}\mathrm{P}_{m}(\mathrm{F}_{1};\mathrm{F}_{0})
\]

这些图表构成 \(\mathbf{P}^{k}(\pi)\) 的一个 \(\mathbf{C}^{r - k}\)-图册，从而赋予 \(\mathbf{P}^{k}(\pi)\) 类 \(\mathbf{C}^{r - k}\) 的向量丛结构，底空间为 X。该向量丛记作 \(\mathbf{P}^{k}(\mathbf{E})\)；其底层流形结构是 12.5.1 中定义的结构。

令 U 为 X 的开子集，且 \(f \in \mathcal{S}_{\mathrm{E}}^{r}(\mathrm{U})\) 是 E 在 U 上的类 \(\mathbf{C}^{r}\) 截面。映射 \(j^{k}(f): x \mapsto j_{x}^{k}(f)\) 是类 \(\mathbf{C}^{r - k}\) 的 \(\mathrm{P}^{k}(\mathrm{E})\) 截面，在 U 上定义。映射 \(j^{k}: \mathcal{S}_{\mathrm{E}}^{r}(\mathrm{U}) \to \mathcal{S}_{\mathrm{P}^{k}(\mathrm{E})}^{r - k}(\mathrm{U})\) 是 K-线性的。

12.6.2. 若 F 是 Banach 空间，则通过 12.5.2 将 \(\mathrm{P}^{k}(\mathrm{F}_{\mathbf{x}})\) 与 \(\mathrm{J}^{k}(\mathrm{X},\mathrm{F})\) 等同，从而赋予后一个流形一个以 X 为底、类为 \(C^{r-k}\) 的向量丛结构。当 F = K 时，用 \(\mathrm{P}^{k}(\mathrm{X})\) 代替 \(\mathrm{P}^{k}(\mathrm{K}_{\mathbf{x}})\)。

例。--- 取 \(X = K^{n}\)，并记 \(u_{1}, \ldots, u_{n}\) 为 \(K^{n}\) 上的坐标函数；于是 \(\mathrm{P}_{0}^{k}(\mathrm{X})\) 是 \(\mathrm{P}^{k}(\mathrm{X})\) 在 0 处的纤维，其一组基由 k 阶喷射组成

单项式 \(u_{1}^{m_1} \ldots u_{n}^{m_n}\) 的喷射，其中 \(m_i \geqslant 0\)，\(\sum_{i=0}^{n} m_i \leqslant k\)。

12.6.3. 设 \(d\) 为满足 \(\geqslant 0\) 的整数，\(\mathrm{E}_{1},\ldots,\mathrm{E}_{d}\) 和 \(\mathbf{F}\) 是类 \(\mathbf{C}^{r}\) 且以 \(\mathbf{X}\) 为底空间的向量丛，并令 \(u\colon\mathrm{E}_{1}\times_{\mathbf{x}}\cdots\times_{\mathbf{x}}\mathrm{E}_{d}\to\mathbf{F}\) 是类 \(\mathbf{C}^{r}\) 的多线性态射（7.3.1）。于是存在唯一的多线性态射

\[
\mathrm{P} ^ {k} (u) \colon \mathrm{P} ^ {k} (\mathrm{E} _ {1}) \times_ {\mathrm{x}} \dots \times_ {\mathrm{x}} \mathrm{P} ^ {k} (\mathrm{E} _ {d}) \rightarrow \mathrm{P} ^ {k} (\mathrm{F})
\]

使得

\[
\mathrm{P} ^ {k} (u) \left(j ^ {k} \left(s _ {1}\right), \dots , j ^ {k} \left(s _ {d}\right)\right) = j ^ {k} \left(u \left(s _ {1}, \dots , s _ {d}\right)\right)
\]

对于 X 的每个开子集 U 以及每列截面 \(s_i \in \mathcal{S}_{\mathbf{E}_i}^r(\mathbf{U})\)（\(1 \leqslant i \leqslant d\)），都有上述等式。

若 A 是以 X 为底空间的代数丛（7.3.2），则由从 \(\mathbf{P}^{k}(\mathbf{A}) \times_{\mathbf{x}} \mathbf{P}^{k}(\mathbf{A})\) 到 \(\mathbf{P}^{k}(\mathbf{A})\) 的态射（由乘法 \(A \times_{x} A \to A\) 导出）使 \(\mathbf{P}^{k}(\mathbf{A})\) 成为代数丛；若 A 是结合代数丛，则 \(\mathbf{P}^{k}(\mathbf{A})\) 是结合代数丛；此外，若 M 是 A-模丛（7.3.3），则 \(\mathbf{P}^{k}(\mathbf{M})\) 是 \(\mathbf{P}^{k}(\mathbf{A})\)-模丛。特别地，\(\mathbf{P}^{k}(\mathbf{X})\) 是结合、交换且有单位元的代数丛；若 E 是向量丛，则 \(\mathbf{P}^{k}(\mathbf{E})\) 是 \(\mathbf{P}^{k}(\mathbf{X})\)-模丛。若 \(E \to F\) 是向量丛的态射，则 \(\mathbf{P}^{k}(u): \mathbf{P}^{k}(\mathbf{E}) \to \mathbf{P}^{k}(\mathbf{F})\) 是 \(\mathbf{P}^{k}(\mathbf{X})\)-同态。

12.6.4. 若 \(E \xrightarrow{u} F \xrightarrow{v} G\) 是以 X 为底的向量丛的局部可裂正合序列，则序列

\[
\mathbf {P} ^ {k} (\mathrm{E}) \rightarrow \mathbf {P} ^ {k} (\mathrm{F}) \rightarrow \mathbf {P} ^ {k} (\mathrm{G})
\]

其中所考虑的同态为 \(\mathrm{P}^{k}(u)\) 和 \(\mathrm{P}^{k}(v)\)。

12.6.5. 令 \(k'\) 和 \(k''\) 为和为 \(k\) 的两个正整数，并令 E 为以 X 为底空间的类 \(C^{r}\) 向量丛。映射

\[
\alpha \colon \mathrm{J} ^ {k} (\mathrm{X}, \mathrm{E}) \rightarrow \mathrm{J} ^ {k ^ {\prime \prime}} (\mathrm{X}, \mathrm{J} ^ {k ^ {\prime}} (\mathrm{X}, \mathrm{E})) \quad (\text { cf. } 1 2. 3. 8)
\]

诱导出一个向量丛态射

\[
\beta \colon \mathbf {P} ^ {k} (\mathrm{E}) \rightarrow \mathbf {P} ^ {k ^ {\prime \prime}} (\mathbf {P} ^ {k ^ {\prime}} (\mathrm{E}))
\]

它属于类 \(C^{r-k}\)。若 U 是 X 的开子集且 \(f\in\mathcal{S}_{\mathrm{E}}^{r}(\mathrm{U})\)，则有

\[
\beta (j ^ {k} (f)) = j ^ {k ^ {\prime \prime}} (j ^ {k ^ {\prime}} (f)).
\]

当 K 的特征为零时，\(\beta\) 是从 \(\mathbf{P}^{k}(\mathbf{E})\) 到 \(\mathbf{P}^{k^{\prime\prime}}(\mathbf{P}^{k^{\prime}}(\mathbf{E}))\) 的一个向量子丛上的同构。

12.6.6. 设 \(k'\) 是满足 \(0 \leqslant k' \leqslant k\) 的整数，E 是类 \(C^r\)、以 X 为底空间的向量丛。映射

\[
r ^ {k, k ^ {\prime}}: \mathrm{P} ^ {k} (\mathrm{E}) \rightarrow \mathrm{P} ^ {k ^ {\prime}} (\mathrm{E})
\]

是类 \(C^{r-k}\) 的局部可裂满射态射。其核 \(\mathbf{N}^{k,k'}(\mathbf{E})\) 由 E 的截面与零截面具有阶数 \(\geqslant k'\) 的接触的喷射组成。

12.6.7（“向量函子 \(P_{m}\)”）。沿用 7.6、7.7 和 7.8 的记号，置 \(I_{+} = \{0\}\)、\(I_{-} = \{1\}\)，并令 m 为满足 \(\geqslant 0\) 的整数。若 \(\mathcal{V} = (\mathrm{V}_{0}, \mathrm{V}_{1})\) 是一对 Banach 空间，记 \(\tau_{m}(\mathcal{V})\) 为 \(\mathrm{P}_{m}(\mathrm{V}_{1}; \mathrm{V}_{0})\)，即定义在 \(V_{1}\) 上、取值于 \(V_{0}\) 的 m 次连续齐次多项式空间（A.2）。同样，若 \(f = (f_{0}, f_{1})\)，其中 \(f_{0}: V_{0} \to V_{0}'\) 和 \(f_{1}: V_{1}' \to V_{1}\) 是 Banach 空间的态射，则记 \(\tau_{m}(f)\) 为映射 \(p \mapsto f_{0} \circ p \circ f_{1}\)，从 \(\mathrm{P}_{m}(\mathrm{V}_{1}; \mathrm{V}_{0})\) 到 \(\mathrm{P}_{m}(\mathrm{V}_{1}', \mathrm{V}_{0}')\)。由此得到向量函子 \(\tau_{m}\)，其类为 \(\mathbf{C}^{\omega}\)。若 \(\mathbf{E}_{0}\) 和 \(\mathbf{E}_{1}\) 是以 \(\mathbf{X}\) 为底空间的两个向量丛，则 \(\mathbf{P}_{m}(\mathbf{E}_{1};\mathbf{E}_{0})\) 是由 \((\mathbf{E}_{0},\mathbf{E}_{1})\) 通过 \(\tau_{m}\) 得到的向量丛（7.6.2）。

12.6.8. 沿用 12.6.1 的记号和假设。存在唯一一个向量丛态射

\[
\iota\colon\mathrm{P}_{k}(\mathrm{T}(\mathrm{X});\mathrm{E})\rightarrow\mathrm{P}^{k}(\mathrm{E})\quad(\text{cf. 12.6.6}),
\]

使得无论 E 的向量图 \(c_{0} = (\mathrm{U}, \varphi, \mathrm{F}_{0})\) 和 X 的图 \(c_{1} = (\mathrm{U}, \psi, \mathrm{F}_{1})\) 如何选取，以下图表都可交换：

\[
\begin{array}{c c c} \mathrm{P} _ {k} (\mathrm{T} (\mathrm{X}), \mathrm{E}) | \mathrm{U} & \xrightarrow {\iota | \mathrm{U}} & \mathrm{P} ^ {k} (\mathrm{E}) | \mathrm{U} \\ \Biggl \downarrow_ {\eta} & & \Biggl \downarrow_ {\theta} \\ \mathrm{U} \times \mathrm{P} _ {k} (\mathrm{F} _ {1}; \mathrm{F} _ {0}) & \xrightarrow {i} & \mathrm{U} \times \prod_ {m = 0} ^ {k} \mathrm{P} _ {m} (\mathrm{F} _ {1}; \mathrm{F} _ {0}) \end{array}
\]

其中：1）\(\iota|_{\mathbf{U}}\) 是 \(\iota\) 在 \(\mathrm{P}_{k}(\mathrm{T}(\mathbf{X}),\mathbf{E})|_{\mathbf{U}}\) 上的限制；

\begin{enumerate}
\def\labelenumi{\arabic{enumi})}
\setcounter{enumi}{1}
\item
  \(\theta\) 是 12.6.1 中定义的双射；
\item
  i 是映射 \((u, p) \mapsto (u, 0, \ldots, 0, p)\)；
\item
  \(\eta\) 是通过向量函子 \(\tau_{k}\) 由 T(X) 的向量图表 \(c_{1}^{\prime}\)（参见 8.1.1）以及 E 的向量图表 \(c_{0}\) 得到的。
\end{enumerate}

映射 \(\iota\) 是从 \(\mathrm{P}_{k}(\mathrm{T}(\mathrm{X});\mathrm{E})\) 到向量子丛 \(\mathbf{N}^{k,k - 1}(\mathbf{E})\)（属于 \(\mathrm{P}^k (\mathbf{E})\)）的同构（参见 12.6.6）；序列

\[
0 \rightarrow \mathrm{P}_{k}(\mathrm{T}(\mathrm{X});\mathrm{E}) \xrightarrow {\iota} \mathrm{P}^{k}(\mathrm{E}) \xrightarrow {r^{k,k-1}} \mathrm{P}^{k-1}(\mathrm{E}) \rightarrow 0
\]

是向量丛的局部可裂正合序列。

更一般地，置 \(\mathrm{N}_{0}=\mathrm{P}^{k}(\mathrm{E})\)，并令 \(\mathrm{N}_{m}=\mathrm{N}^{k,m-1}(\mathrm{E})\) 对 \(m\geqslant1\) 成立，从而 \(N_{m}\) 构成 \(\mathrm{P}^{k}(\mathrm{E})\) 的递减向量子丛序列：

\[
\mathrm{P} ^ {k} (\mathrm{E}) = \mathrm{N} _ {0} \supset \mathrm{N} _ {1} \supset \dots \supset \mathrm{N} _ {k} \supset \mathrm{N} _ {k + 1} = 0.
\]

对于 \(0 \leqslant m \leqslant k\)，投影 \(r^{k,m}: \mathbf{P}^{k}(\mathbf{E}) \to \mathbf{P}^{m}(\mathbf{E})\) 定义了从 \(\mathbf{N}_{m}/\mathbf{N}_{m+1}\) 到 \(\mathbf{N}^{m,m-1}(\mathbf{E})\) 的同构；鉴于上文所述，因而得到同构

\[
\iota_ {m} \colon \mathrm{N} _ {m} / \mathrm{N} _ {m + 1} \rightarrow \mathrm{P} _ {m} (\mathrm{T} (\mathrm{X}); \mathrm{E}) \quad (0 \leqslant m \leqslant k).
\]

特别地，有 \(N_{0}/N_{1} \simeq E\) 和 \(N_{1}/N_{2} \simeq \mathcal{L}(T(X); E)\)。当 k = 1 时，这给出一个正合序列：

\[
0 \rightarrow \mathscr {L} (\mathrm{T} (X); E) \rightarrow \mathrm{P} ^ {1} (E) \rightarrow E \rightarrow 0.
\]

